% English translation of checked_pages/esq_119_015.tex; source: cote 119, site 15, PDF page 16, typescript page 13.
% AI-translated by GPT-6 Astra (OpenAI, ChatGPT Pro mode).
\scanpage{119/15}
\mspage{13}

giving rise to curves of genus zero (which therefore
provide no ``moduli''), one can say that the curve
in question is ``pinned down'' as soon as three of its points are fixed, for
example three vertices of the map, or more generally three
centers of facets (vertices, edges or faces) - from that point on, the structural
map $f : X \to \mathbb{P}^{1}_{\mathbb{C}}$ can be interpreted as a
well-determined rational function
\[ f(z) = P(z)/Q(z) \in \mathbb{C}(z) \]
the quotient of two well-determined relatively prime polynomials,
with $Q$ monic, satisfying algebraic conditions
that express, in particular, the fact that $f$ is
unramified outside the values $0$, $1$, $\infty$, and imply that the
coefficients of these polynomials are \uline{algebraic numbers}; thus
their zeros are algebraic numbers, representing respectively
the vertices and the centers of the faces of the map under
consideration.

Returning to the general case, since finite maps are interpreted
as coverings over $\bar{\mathbb{Q}}$ of an algebraic curve defined
over the prime field $\mathbb{Q}$ itself, it follows that the Galois
group $\mathrm{I}\!\Gamma$ of $\bar{\mathbb{Q}}$ over $\mathbb{Q}$ acts naturally on the category of these
maps. For example, the action of an automorphism
$\gamma \in \mathrm{I}\!\Gamma$ on a spherical map given by the rational function
above is obtained by applying $\gamma$ to the coefficients of the
polynomials $P$, $Q$. Here, then, is this mysterious group $\mathrm{I}\!\Gamma$ intervening as
a transforming agent on topological-combinatorial shapes of the
most elementary kind imaginable, prompting questions
such as: are these given oriented maps ``conjugate'',
or: what exactly are the conjugates\frapparatus{V-FR-015-001} of this given oriented
map ? (There are, visibly, only finitely many.)

I have dealt with a few concrete cases (for coverings
of low degrees) by various expedients; J. Malgoire has dealt with
a few others - I doubt that there is a uniform method of
answering them by computer calculations. My reflection very quickly

\margin{15}

